$G=(V,E)$ mit $V$: Menge der Knoten ($n$), $G$: Menge der Kanten ($m$).
\begin{sectionbox}
    \begin{description}
        \item[v adjazent zu u] $(u,v)$ bzw. $\{u,v\} \in E$
        \item[inzident] Knoten $v$ und Kante $e$, wenn $e=\{x,v\}$
    \end{description}
\end{sectionbox}

\subsection{Darstellung von Graphen}
\begin{sectionbox}
    \subsubsection{Adjazenzmatrix}
    Speichert vorhandenen Kanten in $A \in \mathbb R^{n\times n}$ Matrix. Eine Zeile gehört zu einem Knoten
    \begin{itemize}
        \item $A(i,j)=1(=\text{Gewicht, bei gewichtetem Graph})$, falls Kante von $v_i$ zu $v_j$ existiert
        \item $A(i,j)=0$, sonst
    \end{itemize}
    Symmetrisch bei ungerichtetem Graphen. Kanten einfügen löschen einfach, Knoten einfügen/löschen schwierig. Hohen Speicherverbrauch

    \subsubsection{Adjazenzliste}
    Speichert für jeden Knoten eine verkette Liste mit Nachbarn des jeweiligen Knotens; die Knoten selbst werden woanders gespeichert (Array z.B.) \\

    \subsubsection{Statischer Graph}
    Array mit allen Kanten, wobei ein Bereich im Array zu den Verbindungen eines bestimmten Knotens gehört. \\
    Array mit allen Knoten, wobei für jeden Knoten der Index seines Bereichs und die Menge an Nachbarn gespeichert wird
\end{sectionbox}
